\documentclass[../../main.tex]{subfiles}
\begin{document}

{\bf [解]}

題意の条件から以下
\begin{align}
\dfrac{1}{p} + \dfrac{1}{q} = 1, \, 0 < p, \, 0 < q \label{eq:1}
\end{align}
で考える．この領域を下図に示す．

\begin{figure}[htb]
\begin{center}
\begin{tikzpicture}[scale=1.5, >=stealth]
  % -------------------------------------------------------------
  % 座標軸
  % -------------------------------------------------------------
  \draw[->] (-0.3, 0) -- (4.2, 0) node[right] {$p$};
  \draw[->] (0, -0.3) -- (0, 4.2) node[above] {$q$};
  \node[below left, font=\small] at (0,0) {$O$};

  % -------------------------------------------------------------
  % 漸近線 (p = 1, q = 1)
  % -------------------------------------------------------------
  \draw[dashed, gray!80] (1, 0) node[below, font=\small, black] {$1$} -- (1, 4.0);
  \draw[dashed, gray!80] (0, 1) node[left, font=\small, black] {$1$} -- (4.0, 1);

  % -------------------------------------------------------------
  % 双曲線 1/p + 1/q = 1  <=>  q = p / (p - 1)  (p > 1 のブランチ)
  % -------------------------------------------------------------
  \draw[thick, red!80!black, domain=1.28:4.0, samples=100]
      plot (\x, {\x / (\x - 1)}) 
      node[right, font=\small] {$\dfrac{1}{p} + \dfrac{1}{q} = 1$};

  % -------------------------------------------------------------
  % 特徴的な点 (例: p=2, q=2 の対称点) のプロットと補助線
  % -------------------------------------------------------------
  \coordinate (SymmetricPoint) at (2, 2);
  \fill[purple] (SymmetricPoint) circle (1.5pt);

  \draw[dotted, gray!70] (2, 0) node[below, font=\small, black] {$2$} -- (SymmetricPoint);
  \draw[dotted, gray!70] (0, 2) node[left, font=\small, black] {$2$} -- (SymmetricPoint);

  % -------------------------------------------------------------
  % 1 < p, 1 < q の領域であることを示す強調マーク・補足
  % -------------------------------------------------------------
  % \node[font=\small, purple!80!black, fill=white, inner sep=2pt] 
  %    at (2.7, 2.7) {$1 < p$ かつ $1 < q$};

\end{tikzpicture}
\end{center}
\caption{\cref{eq:1}の描く領域}
\end{figure}

簡単のため
\begin{align}
 f(x) = \dfrac{1}{p}x^p - x + \dfrac{1}{q}
\end{align}
とおくと，この正負を考えればよい．
一階微分は
\begin{align}
 f'(x) = x^{p-1} - 1
\end{align}
であり，\cref{eq:1}より$1<p$だから増減表は以下のようになる．

\begin{table}[htb]
\centering
\begin{tabular}{c|c*3{c}}
$x$ & $0$ & $\dots$ & $1$ & $\dots$ \\ \hline
$f'$ & & $-$ & $0$ & $+$ \\ \hline
$f$ & & $\searrow$ & $0$ & $\nearrow$
\end{tabular}
\end{table}

よって$f(x)\ge 0$であり，等号成立は$x=1$の時だから求める大小関係は
\begin{align}
\begin{cases}
  x = \dfrac{1}{p}x^p + \dfrac{1}{q} & \left(x = 1 \text{ の時 }\right) \\[10pt]
  x < \dfrac{1}{p}x^p + \dfrac{1}{q} & \left(x \ne 1 \text{ の時 }\right)
\end{cases}
\end{align}
である．

\end{document}
